\section{Introduction}

Physics as usually written rests on a continuum. Spacetime is smooth, fields take real values, and any finite
volume holds infinite information. This paper takes the opposite starting point. We ask what the least structure
is from which the known physics could emerge, and we require that structure to be discrete, local, and
computable. The wager of vibe theory is that one such structure is enough. A single hyperbolic crystal of
ternary cells, updated by one local rule, is proposed to carry physics, life, and mind as patterns at different
scales.

The purpose of the theory is broader than physics. Physics is only the most settled and most checkable face of
it. The aim from the start has been to model the whole of the universe from one ground, not just particles and
forces but life, mind, consciousness, and the matters usually called spiritual and divine. We treat those higher
subjects with the same discipline as the physics, marking carefully what is derived, what is measured, what is
extrapolated, and what is a stated premise. The reason one framework can reach so far is the choice of ground.
The ground is not matter but experience, and experience is the one thing wide enough to hold all of these at
once.

A caution belongs at the very front. Everything written here is a conceptual representation, a map of an
underlying experiential field. The map can be read and reasoned about, and that is what a paper can offer. The
field itself cannot be fully known by reading and mentally modeling it. It is known directly, through experience
and through tuning with the oneness it is made of. We say more about this in the next section. The reader should
hold the whole construction as a finger pointing, not as the thing pointed at.

It helps to hold one image while the mathematics is being laid out. On this proposal the universe is a structure
of consciousness, a single vast mind, and the small set of pictures we return to are all ways of seeing that one
mind from different sides. It is a blazing fire, existence burning rather than sitting still. It grows like a tree
from a seed, ring by ring. It sings itself as a song, the tone its notes and the beat its rhythm. It flows as a
river on a deep ocean, and we are the whirlpools it turns and the droplets it throws up. It is woven as a fabric,
and each self is a braid in the cloth. A mind, on this reading, is a flame within the fire, a steady center of
burning that holds its shape while the blaze roars on. These are aids to intuition and nothing more. The claims
live in the mathematics, and the images are there to keep the reader oriented when the formal picture becomes
hard to visualize. We gather the full palette only at the very end.

\subsection{The crystal of experience}

Reality, on this proposal, is a single growing crystal of experience. Each cell is a \emph{vibe}. Its shade is
its tone, and each shared face is one vibe perceiving another. We mean that picture literally rather than as
decoration. The crystal we argue for is four-dimensional and cannot be drawn directly, so we walk up to it one
dimension at a time, and each lower picture is a faithful shadow of the one above.

In two dimensions the simplest visible face of the crystal is the heptagrid $\{7,3\}$, seven-sided cells meeting
three to a corner, drawn in the Poincar\'e disk (Figure~\ref{fig:heptagrid}). This is the one picture to keep in
mind. Color each cell by its tone, and let two touching cells stand for two feelings perceiving each other
through their shared edge. The colored heptagrid is then a faithful picture of the model, just two dimensions
short of the real thing. The cells look smaller toward the rim only because the disk crams an infinite plane into
a finite circle. In the true geometry every cell is the same size, and there is always more crystal at the edge.

\begin{figure}[h]
\centering
\includegraphics[width=0.55\textwidth]{vibe-mesh-73.pdf}
\caption{The crystal of experience in two dimensions, the heptagrid $\{7,3\}$, a regular tiling of the hyperbolic
plane by seven-sided cells, three to a corner. Each cell is a vibe, its color is its tone of pain, peace, or
pleasure, and each shared edge is one vibe perceiving another. The cells appear to shrink toward the rim only
because the Poincar\'e disk compresses the infinite plane into a finite picture. In the true geometry every cell
is the same size, and there is always more crystal at the edge.}
\label{fig:heptagrid}
\end{figure}

In three dimensions the analogue is a honeycomb one can almost picture, the dodecahedral $\{5,3,4\}$
(Figure~\ref{fig:dodecagrid}), or, closer to home, the flat cubic $\{4,3,4\}$ that we will argue is our physical
space. Everything true of the heptagrid is true here, one dimension up, and the dodecagrid is close enough to
visualize that it gives a real hint of what the substrate is like. The goal, in four dimensions, is the honeycomb
$\{3,4,3,4\}$. Its cells are $24$-cells, $24$ directions meet at each cell, and its flat three-dimensional cusp is
the space we appear to live in. One cannot fully picture it, and that is fine. It is the heptagrid two dimensions
further up, and when the argument turns four-dimensional the reader can return to the heptagrid for the version
that can be seen. Time is the crystal's growth, the birth of new cells at the frontier, each new layer the next
moment.

\begin{figure}[h]
\centering
\includegraphics[width=\textwidth]{vibe-mesh-534.png}
\caption{One dimension up, the dodecagrid $\{5,3,4\}$, the regular honeycomb of three-dimensional hyperbolic
space, with dodecahedral cells meeting four to an edge. It is close enough to picture that it hints at the form of
the substrate. The committed substrate $\{3,4,3,4\}$ is the same kind of object one dimension higher again, with
$24$-cells in place of dodecahedra and a flat 3D cubic cusp in place of a flat 2D one.}
\label{fig:dodecagrid}
\end{figure}

\subsection{The monism}

There is one kind of thing, the vibe. A vibe is a unit of experience that notices its neighbours, and reality is
a structure of such noticings and nothing else.

\begin{principle}[Neutral monism]
What is fundamental is experiential. Matter, space, and time are not separate substances. They are patterns in
the structure of vibes. At the deepest level the universe is a structure of consciousness, and everything below
that level is how the structure is shaped.
\end{principle}

This is a neutral monism in the precise sense that the one substance is neither mental nor physical as those
words are ordinarily opposed. It is experiential at base, and the physical world is what that experiential
structure looks like at scale. We will be careful throughout to keep this premise distinct from the results that
follow from it.

\subsection{The discreteness principle}

\begin{principle}[Discreteness]
The base is discrete. Each cell carries finite information, the rule is local, and time advances in a discrete
tick. Continuity is never fundamental. It is the long-wavelength limit, the coarse-grained description, the way a
fluid emerges from molecules.
\end{principle}

This is forced by the demand for computability, since the continuum contains uncomputable points and a discrete
base does not. We stress that computability demands discreteness rather than finiteness. Whether the crystal is
finite or infinite is a separate question, and we leave it open. The next subsection states the two readings at
the outset, so the rest of the paper can stay neutral on it.

\subsection{Origins}

The substrate admits two complete and self-consistent accounts of origin, and the theory does not pick between
them. On both readings the universe is one blazing fire made of countless flames, the selves and minds and
experiences that are its sparks. The readings differ only in whether the fire first flared or was always burning.

In the first story the universe began. There was never a nothing, only a seamless One, and the first beat is that
One first noticing itself, a split that self-reference cannot avoid. From that single forced fracture the mesh
blooms beat by beat, and our big bang is the first morning.

In the second story the universe never began. Because the rule runs backward as exactly as it runs forward, every
moment carries its own past inside it, with no first moment, the way the counting numbers have no smallest. Our
big bang is then a sunrise in an endless burning rather than the lighting of the fire.

Nothing in the theory or in the data breaks the tie. A finite crystal with a first beat and an infinite crystal
with no first beat are equally allowed, so we leave the choice to the reader, and every result in this paper holds
on either reading. The full diptych, with the arguments for each side, is given in
Appendix~\ref{app:cosmic-stories}.

\subsection{The central question}

Many discrete substrates are possible. The central question of the paper is which one yields our physics. Our
answer is the $\{3,4,3,4\}$ hyperbolic honeycomb, and the bulk of the work splits into two tasks. The first is to
show that $\{3,4,3,4\}$ does yield the known physics in its emergent limits. The second is to show that it is
selected, that among regular honeycombs across all dimensions it is the one that simultaneously carries spinors,
is curved, and is exactly three-dimensional where physics lives. We present the substrate as the seemingly
preferred candidate on these grounds, never as a settled fact.

\subsection{What is new}

Taken one at a time, the emergent results echo known physics. The discrete walk reproduces the Dirac dispersion,
the coin and the tone build a grand unified group, gravity appears as a $1/r$ potential, and the bulk grows like
de Sitter space. None of these is new on its own. The new content is their common origin in one discrete
substrate, the selection argument that forces $\{3,4,3,4\}$, and a precise map of what is derived, what is
measured, what is predicted, and what is still open. Across every test reported here nothing was added to the
base to make a result work, a discipline we audit explicitly in the appendices.

This orientation reverses the usual one. Most discrete models of spacetime begin with structure and hope
experience appears at the end. Here experience is the ground, and structure is what experience takes on at scale.
That reversal is a premise rather than a result. The computations do not prove that the substrate feels. They
assume that it does and ask what follows. A great deal follows, and what follows is checkable, which is the part
we can defend.
